% Part: first-order-logic
% Chapter: axiomatic-deduction
% Section: provability-consistency

\documentclass[../../../include/open-logic-section]{subfiles}

\begin{document}

\iftag{FOL}
      {\olfileid[hi]{fol}{axd}{prv}}
      {\olfileid[hi]{pl}{axd}{prv}}

\olsection{\usetoken{S}{derivability} और संगतता}

अब !!{derivability} संबंध के अनेक गुण स्थापित करेंगे। वे अपने-आप में रोचक
हैं, पर उनमें से प्रत्येक पूर्णता प्रमेय के प्रमाण में भूमिका निभाएगा।

\begin{prop}\ollabel{prop:provability-contr}
  यदि $\Gamma \Proves !A$ और $\Gamma \cup \{!A\}$ असंगत है, तो
  $\Gamma$ असंगत है।
\end{prop}

\begin{proof}
  यदि $\Gamma \cup \{!A\}$ असंगत है, तो $\Gamma \cup \{!A\}
  \Proves \lfalse$। \olref[ptn]{prop:reflexivity} से $\Gamma \Proves
  !B$, प्रत्येक $!B \in \Gamma$ के लिए। परिकल्पना से $\Gamma \Proves !A$
  भी है, इसलिए $\Gamma \Proves !B$, प्रत्येक $!B \in \Gamma \cup
  \{!A\}$ के लिए। \olref[ptn]{prop:transitivity} से $\Gamma \Proves
  \lfalse$; अर्थात् $\Gamma$ असंगत है।
\end{proof}

\begin{prop}
\ollabel{prop:prov-incons}
$\Gamma \Proves !A$ तभी और केवल तभी जब $\Gamma \cup \{\lnot !A\}$
असंगत हो।
\end{prop}

\begin{proof}
पहले मान लें $\Gamma \Proves !A$। तब \olref[ptn]{prop:monotonicity} से
$\Gamma \cup \{\lnot !A\}
\Proves !A$। \olref[ptn]{prop:reflexivity} से $\Gamma \cup \{\lnot
!A\} \Proves \lnot !A$। \olref[prp]{ax:lnot2} से हमारे पास
$\Proves \lnot !A \lif (!A \lif \lfalse)$ भी है। अतः
\olref[ded]{prop:mp} के दो प्रयोगों से $\Gamma \cup \{\lnot
!A\} \Proves \lfalse$।

अब मान लें $\Gamma \cup \{\lnot !A\}$ असंगत है, अर्थात् $\Gamma
\cup \{\lnot !A\} \Proves \lfalse$। निगमन प्रमेय से $\Gamma
\Proves \lnot !A \lif \lfalse$। \olref[prp]{ax:lfalse2} से
$\Gamma \Proves (\lnot !A \lif
\lfalse) \lif \lnot\lnot !A$, इसलिए \olref[ded]{prop:mp} से $\Gamma
\Proves \lnot\lnot !A$। चूँकि $\Gamma \Proves
\lnot\lnot !A \lif !A$ (\olref[prp]{ax:dne}), इसलिए
\olref[ded]{prop:mp} के फिर प्रयोग से $\Gamma
\Proves !A$।
\end{proof}

\begin{prob}
सिद्ध कीजिए कि $\Gamma \Proves \lnot !A$ तभी और केवल तभी जब
$\Gamma \cup \{!A\}$ असंगत हो।
\end{prob}

\begin{prop}\ollabel{prop:explicit-inc}
  यदि $\Gamma \Proves !A$ और $\lnot !A \in \Gamma$, तो $\Gamma$ असंगत है।
\end{prop}

\begin{proof}
  \olref[prp]{ax:lnot2} से
  $\Gamma \Proves \lnot !A \lif (!A \lif \lfalse)$।
  \olref[ded]{prop:mp} के दो प्रयोगों से $\Gamma \Proves \lfalse$।
\end{proof}


\begin{prop}\ollabel{prop:provability-exhaustive}
  यदि $\Gamma \cup \{!A\}$ और $\Gamma \cup \{\lnot !A\}$ दोनों असंगत
  हैं, तो $\Gamma$ असंगत है।
\end{prop}

\begin{proof}
अभ्यास।
\end{proof}

\tagprob{FOL}
\begin{prob}
  \olref[fol][axd][prv]{prop:provability-exhaustive} सिद्ध कीजिए।
\end{prob}
\tagendprob

\tagprob{notFOL}
\begin{prob}
  \olref[pl][axd][prv]{prop:provability-exhaustive} सिद्ध कीजिए।
\end{prob}
\tagendprob


\end{document}
